\documentclass[aspectratio=169,10pt]{beamer}
\newcommand{\reporoot}{../..}
% Shared preamble of the main deck and the topic decks.
% Define \reporoot (the path from the .tex file to the repository root)
% before \input-ing this file.
\usetheme[progressbar=frametitle,numbering=fraction,sectionpage=progressbar]{metropolis}

\usepackage{amsmath,amssymb}
\usepackage{booktabs}
\usepackage{array}
\newcolumntype{L}[1]{>{\raggedright\arraybackslash}p{#1}}
\usepackage{graphicx}
\usepackage{tikz}
\usetikzlibrary{arrows.meta,positioning,decorations.pathreplacing}
\usepackage[authoryear,round]{natbib}
\usepackage{xcolor}

\graphicspath{{\reporoot/figures/}}

% Every simulated number on the slides is produced by experiments/*.py and read from here.
\input{\reporoot/results/e01_estimators.tex}
\input{\reporoot/results/e02_scaling.tex}
\input{\reporoot/results/e03_var_backtests.tex}
\input{\reporoot/results/e04_es_backtest.tex}
\input{\reporoot/results/e05_pla.tex}
\input{\reporoot/results/e06_estimation_risk.tex}
\input{\reporoot/results/e07_model_uncertainty.tex}
\input{\reporoot/results/e08_asrf.tex}
\input{\reporoot/results/e09_copula_is.tex}
\input{\reporoot/results/e10_calibration_tests.tex}
\input{\reporoot/results/e11_discrimination.tex}
\input{\reporoot/results/e12_exposure.tex}
\input{\reporoot/results/e13_cva_wwr.tex}
\input{\reporoot/results/e14_nested.tex}
\input{\reporoot/results/e15_variance_reduction.tex}
\input{\reporoot/results/e16_lda.tex}
\input{\reporoot/results/e17_evt.tex}
\input{\reporoot/results/e18_reverse_stress.tex}
\input{\reporoot/results/e19_regimes.tex}
\input{\reporoot/results/gallery.tex}

\newcommand{\E}{\mathbb{E}}
\newcommand{\Var}{\operatorname{Var}}
\newcommand{\VaR}{\mathrm{VaR}}
\newcommand{\ES}{\mathrm{ES}}
\newcommand{\dd}{\mathrm{d}}
\definecolor{accent}{HTML}{0C5DA5}
\definecolor{warn}{HTML}{C0392B}
\newcommand{\hl}[1]{\textcolor{accent}{\textbf{#1}}}

\setbeamerfont{footnote}{size=\tiny}
\setbeamerfont{caption}{size=\scriptsize}

\author{Yaohui Shu}
\institute{Stony Brook University}
\date{November 2026}

% Topic decks: a references frame that lists only the works cited in the deck.
\newcommand{\deckreferences}{%
  \begin{frame}[allowframebreaks]{References}
    \tiny
    \renewcommand{\bibsection}{}
    \bibliographystyle{plainnat}
    \bibliography{\reporoot/slides/refs}
  \end{frame}}

% Topic decks: \decknumber{n} puts the deck number on the title page.
\newcommand{\decknumber}[1]{\institute{Stony Brook University\hfill{\scriptsize\color{gray}Deep dive #1 of 9}}}


\title{Operational Risk and Extreme Values}
\subtitle{The loss distribution approach, the single-loss approximation, and tail estimation}
\decknumber{8}

\begin{document}

\maketitle

\begin{frame}{The question}
  \small
  \begin{columns}[T]
    \begin{column}{0.55\textwidth}
      Fraud, outages, legal settlements: a few huge losses dominate the total. \hl{How do we estimate a one-in-a-thousand-year annual loss from a few hundred recorded events, and how sure can we be?}
      \begin{itemize}
        \item Heavy tails make the largest loss the main driver of the annual total.
        \item Extreme value theory tells us how to extrapolate beyond the data, and how uncertain that is.
      \end{itemize}
    \end{column}
    \begin{column}{0.42\textwidth}
      \begin{alertblock}{In this deck}
        \small
        \begin{enumerate}
          \item The loss distribution approach.
          \item Exact compound distributions: FFT and Panjer.
          \item The single-loss approximation.
          \item Capital estimation risk.
          \item Peaks over threshold and the Hill estimator.
        \end{enumerate}
      \end{alertblock}
    \end{column}
  \end{columns}
\end{frame}

\begin{frame}{Operational losses: one year, and a million years}
  \centering
  \includegraphics[width=0.95\textwidth,height=0.54\textheight,keepaspectratio]{gallery/g08_operational_risk.png}\par\vspace{0.4em}
  \footnotesize\raggedright
  (a) One simulated year with 25 losses a year on average and lognormal sizes (log-mean 10, log-standard deviation $\sigma = 2$): the largest of its \GalEightEvents{} losses is \GalEightShare\% of the year's total. Grey: 39 other years. (b) Annual loss distributions for $\sigma$ from 1.0 to 2.6, one million simulated years each, on log scales; red is the tail beyond the 99.9\% quantile.
\end{frame}

\begin{frame}{The loss distribution approach}
  \small
  Annual loss $S = X_1 + \dots + X_N$, with frequency $N \sim \mathrm{Poisson}(\lambda)$ (events per year) and i.i.d.\ severities $X_k$ (currency) with distribution function $F$, independent of $N$.
  \begin{columns}[T]
    \begin{column}{0.5\textwidth}
      \textbf{Exact distribution.} Discretize $F$ on a grid of span $h$ and compute the compound Poisson law by the fast Fourier transform,
      \[ f_S = \mathrm{IFFT}\big(\exp(\lambda(\mathrm{FFT}(f_X) - 1))\big), \]
      with zero padding against wrap-around. Panjer's recursion
      \[ f_S(k) = \frac{\lambda}{k}\sum_{j=1}^{k} j\,f_X(j)\,f_S(k-j) \]
      gives the same numbers (tested to $10^{-10}$).
    \end{column}
    \begin{column}{0.46\textwidth}
      \textbf{Single-loss approximation} \citep{bocker2005operational}. For subexponential (heavy-tailed) severities,
      \[ \VaR_\alpha(S) \approx F^{-1}\!\left(1 - \frac{1-\alpha}{\lambda}\right) \quad (\alpha \to 1): \]
      at very high levels the annual total is large because one loss is. The capital is a severity quantile, so its accuracy hinges on the severity tail.
    \end{column}
  \end{columns}
\end{frame}

\begin{frame}{How good is the approximation, and how stable is the capital?}
  \centering
  \includegraphics[width=0.86\textwidth]{fig16_lda.png}
  \begin{columns}[T]
    \begin{column}{0.5\textwidth}
      \footnotesize
      (a) $\lambda = 25$ a year, lognormal severities with $\mu = 10$. At 99.9\% the approximation over the exact value is $\ESixteenSlaOneFive$ ($\sigma = 1.5$), $\ESixteenSlaTwoNine$ ($\sigma = 2$) and $\ESixteenSlaTwoFive$ ($\sigma = 2.5$); at 99\% and $\sigma = 2$, $\ESixteenSlaTwoTwo$. It improves with the level and with tail heaviness, and errs on the low side.
    \end{column}
    \begin{column}{0.46\textwidth}
      \footnotesize
      (b) True model $\sigma = 2$; the bank fits a lognormal to $n$ observed losses and computes 99.9\% capital by the SLA ($\ESixteenTrue$ million at the true parameters). The 90\% range of the estimate over that value is $\ESixteenLoHundred$ to $\ESixteenHiHundred$ with $n = 100$ and $\ESixteenLoThousand$ to $\ESixteenHiThousand$ with $n = 1{,}000$, with the model correctly specified.
    \end{column}
  \end{columns}
\end{frame}

\begin{frame}{Extreme value theory}
  \small
  \begin{columns}[T]
    \begin{column}{0.52\textwidth}
      \textbf{Peaks over threshold.} Excesses $Y = L - u$ over a high threshold $u$ are approximately generalized Pareto,
      \[ P(Y \le y) = 1 - (1 + \xi y/\beta)^{-1/\xi}, \]
      shape $\xi$ (tail index; $\xi > 0$ for heavy tails), scale $\beta$ (loss units). With $N_u$ of $n$ losses above $u$,
      \[ \VaR_\alpha = u + \frac{\beta}{\xi}\left[\left(\frac{n(1-\alpha)}{N_u}\right)^{-\xi} - 1\right] \]
      \citep{mcneil2015quantitative}.
    \end{column}
    \begin{column}{0.45\textwidth}
      \textbf{Hill estimator} from the $k$ largest losses,
      \[ \hat\xi = \frac1k\sum_{i=1}^{k}\ln X_{(n-i+1)} - \ln X_{(n-k)} , \]
      for Pareto-type tails.
      \vspace{0.3em}

      \textbf{The trade-off.} A high threshold (small $k$) keeps the model accurate but leaves few points; a low one adds points the model does not describe. There is no automatic answer.
    \end{column}
  \end{columns}
\end{frame}

\begin{frame}{Estimating a 99.9\% quantile}
  \centering
  \includegraphics[width=0.86\textwidth]{fig17_evt.png}
  \begin{columns}[T]
    \begin{column}{0.5\textwidth}
      \footnotesize
      (a) Student $t_3$ losses ($\xi = 1/3$), 1,000 samples per size. RMSE of $\ln(\hat\VaR/\VaR)$ at $n = 1{,}000$: empirical $\ESeventeenEmpThousand$, normal fit $\ESeventeenNormThousand$, POT at the 90\% threshold $\ESeventeenPotThousand$. At $n = 250$: empirical $\ESeventeenEmpQuarter$, POT $\ESeventeenPotQuarter$.
    \end{column}
    \begin{column}{0.46\textwidth}
      \footnotesize
      (b) One sample of 2,500: the Hill estimate moves from $\ESeventeenHillSmall$ (30 points) to $\ESeventeenHillLarge$ (1,000 points), against the true 0.33. Using too much of the sample biases $\hat\xi$, because the body of a $t$ distribution is not Pareto.
    \end{column}
  \end{columns}
\end{frame}

\begin{frame}{Discussion and takeaways}
  \small
  \begin{columns}[T]
    \begin{column}{0.5\textwidth}
      \textbf{What the experiments show}
      \begin{itemize}
        \item The single-loss approximation is a good capital proxy at 99.9\% for heavy tails, and understates at lower levels.
        \item Even a correctly specified severity model gives 99.9\% capital that varies by a factor of several across samples of realistic size.
        \item A normal model is wrong by a large factor in heavy tails; EVT helps but needs a threshold choice.
      \end{itemize}
    \end{column}
    \begin{column}{0.46\textwidth}
      \textbf{For an operational risk modeler}
      \begin{enumerate}
        \item Report capital with an interval, from bootstrap or from the estimator's sampling distribution.
        \item Show threshold sensitivity (Hill or GPD-shape plots) in every EVT fit.
        \item Use the exact FFT distribution to check approximations.
      \end{enumerate}
      \textbf{Open question.} Combining internal losses, external data and scenarios into one tail estimate with honest uncertainty.
    \end{column}
  \end{columns}
\end{frame}

\begin{frame}[fragile]{Reproduce}
  \small
  Code: \texttt{oprisk.py}, \texttt{evt.py}. Experiments: \texttt{e16\_lda.py}, \texttt{e17\_evt.py}.
\begin{verbatim}
from stochastic_risk import oprisk as op, evt
fx = op.discretize(severity_cdf, h, n)
fs = op.compound_poisson_fft(lam=25, fx=fx)
op.sla_var(0.999, lam=25, severity_ppf=ppf)
evt.pot_var_es(losses, alpha=0.999, threshold_quantile=0.9)
evt.hill(losses, k=200)
\end{verbatim}
\end{frame}

\deckreferences
\end{document}
